%% language-es-test — Spanish multilingual paper (acmart `language=spanish`).
%% Spanish main language with an English translated title/keywords. Verifies the
%% Spanish fixed strings — \keywordsname ("Palabras y Frases Claves
%% Adicionales"), \proofname ("Demostración"), \acksname ("Expresiones de
%% gratitud"), \tablename ("Cuadro") — while the figure name stays "Fig."
%% (acmart sets it globally, acmart.dtx:4199). Twin of language-es-test.typ.
\documentclass[acmsmall,language=spanish]{acmart}

\acmJournal{JACM}
\acmVolume{37}
\acmNumber{4}
\acmArticle{111}
\acmYear{2018}
\acmMonth{8}
\acmDOI{XXXXXXX.XXXXXXX}
\setcopyright{acmlicensed}
\copyrightyear{2018}

\begin{document}

\title{Una nota sobre la complejidad computacional}
\translatedtitle{english}{A note on computational complexity}

\author{Juan Pérez}
\email{perez@example.es}
\affiliation{%
  \institution{Instituto de Investigación}
  \city{Madrid}
  \country{Spain}}

\renewcommand{\shortauthors}{Pérez}

\begin{abstract}
  Una breve nota en español para ejercitar el modo multilingüe de la clase acmart
  con cadenas de caracteres traducidas.
\end{abstract}

\keywords{complejidad, algoritmos, cómputo}
\translatedkeywords{english}{complexity, algorithms, computation}

\maketitle

\section{Introducción}
Este documento comprueba que las cadenas localizadas coinciden entre la clase
LaTeX y el porte de Typst.

\begin{figure}[h]
\centering
\rule{4cm}{2cm}
\caption{Una figura; la etiqueta permanece en inglés.}
\end{figure}

\begin{table}[h]
\caption{Un cuadro con etiqueta localizada.}
\centering
\begin{tabular}{ll}
\toprule
A & B\\
\midrule
1 & 2\\
\bottomrule
\end{tabular}
\end{table}

\begin{proof}
Una breve demostración que termina con un cuadrado CQD.
\end{proof}

\begin{acks}
Agradecemos a los revisores.
\end{acks}

\end{document}
